\centering
\renewcommand{\arraystretch}{1.4}
\begin{tabular}{|l|p{4cm}|p{7.5cm}|}
\hline
\textbf{Tipo de mensaje} & \textbf{Campos} & \textbf{Descripción} \\
\hline
\textbf{Petición} & \texttt{cmd} (+ params) &
  Cliente \textrightarrow\ servidor. \texttt{cmd} identifica la operación; los parámetros opcionales se incluyen en el mismo objeto. \\
\hline
\textbf{Respuesta OK} & \texttt{status="ok"}, \texttt{data} &
  Servidor \textrightarrow\ cliente. \texttt{data} contiene el resultado (escalar, objeto o lista). \\
\hline
\textbf{Respuesta error} & \texttt{status="error"}, \texttt{message} &
  Servidor \textrightarrow\ cliente. \texttt{message} describe la causa del fallo. \\
\hline
\textbf{Push} & \texttt{type} (+ datos) &
  Servidor \textrightarrow\ cliente sin petición previa (\texttt{samples}, \texttt{error}, \texttt{heartbeat}). \\
\hline
\end{tabular}
\vspace{0.6em}

\begin{tabular}{lp{12cm}}
\textbf{Ejemplo petición:}  & \texttt{\{"cmd":"uds\_read\_did","did":"0xF190"\}\textbackslash n} \\
\textbf{Ejemplo respuesta:} & \texttt{\{"status":"ok","data":\{"did":"0xF190","name":"VIN",...\}\}\textbackslash n} \\
\end{tabular}
\vspace{0.5em}
